
% =====================================================================
\section{Individuation by negation, and its non-instantaneity}
\label{sec:negation}
% =====================================================================

Before a separator can have thickness there must be something to separate. This
section settles how a part comes to be at all in a whole with no distinguished
element, and proves two facts that drive the rest of the paper: individuation is
forced to proceed by negation (\Cref{thm:negation-unique}), and no individuation
that yields a genuine part can be completed in a single step
(\Cref{thm:noninstant}), so that every individuation leaves an untouched
residue.

\subsection{Individuation without a selector}

\begin{definition}[Individuation]
\label{def:individuation}
An \emph{individuation} of a part from the whole $\Space$ is a determination of a
region $A$ together with a separator certifying $A$ as distinct from
$\Space\setminus A$. A \emph{selector} is a rule that fixes $A$ by naming an
element or distinguished point of $A$ in advance of the separation.
\end{definition}

\begin{definition}[Negation-sequence]
\label{def:negation}
For a region $A\subseteq\Space$, its \emph{negation-sequence} is its complement
together with the separator that certifies the complement,
\[
\Neg(A)\;:=\;\Space\setminus A,
\]
the totality of what $A$ is not. We say $A$ is \emph{individuated by negation}
if $A$ is determined as the unique region whose negation-sequence is $\Neg(A)$;
that is, $A=\Space\setminus\Neg(A)$, with $\Neg(A)$ given first.
\end{definition}

The point of \Cref{def:negation} is that $\Neg(A)$ requires no selector: it is
``everything else,'' which is completely specified by the whole and the act of
withholding $A$, with no element singled out. We now show this is the only
selector-free route.

\begin{theorem}[Negation is the unique selector-free individuation]
\label{thm:negation-unique}
Let $\Space$ contain no distinguished element \textup{(}no point, region, or
order is given in advance\textup{)}. Then any individuation of a part $A$ that
uses no selector determines $A$ as $\Space\setminus\Neg(A)$ for its
negation-sequence $\Neg(A)$; and conversely every $\Neg(A)$ determines a unique
such $A$. No selector-based rule individuates a part without invoking a further
selector.
\end{theorem}

\begin{proof}
\emph{(Necessity.)} Suppose an individuation fixes $A$ without a selector. To
fix $A$ is to decide, of each part of $\Space$, whether it belongs to $A$ or to
the remainder. A positive rule---one that fixes $A$ by asserting membership of
some element or distinguished feature---names that element or feature; by
hypothesis no element or feature is distinguished in advance, so the rule must
itself supply one. Supplying a distinguished element is exactly providing a
selector, and that selector must in turn be fixed: either it is given in advance
(contradicting that none is) or it is fixed by a prior rule, which faces the same
dichotomy. The regress terminates only if some rule fixes $A$ without naming any
element of $A$. The sole such rule available is the withholding of $A$ from the
whole: specifying the totality $\Space\setminus A=\Neg(A)$ that $A$ is to be
distinct from, and taking $A$ as its complement. This names no element of $A$;
it names only ``the rest,'' which the whole already supplies. Hence
$A=\Space\setminus\Neg(A)$.

\emph{(Sufficiency and uniqueness.)} Given a negation-sequence
$\Neg(A)\subsetneq\Space$ with $\mu(\Space\setminus\Neg(A))>0$, the region
$A:=\Space\setminus\Neg(A)$ is determined uniquely as a set, since complementation
in $\Space$ is an involution: $\Space\setminus(\Space\setminus\Neg(A))=\Neg(A)$.
Thus $\Neg$ is a bijection between admissible parts and admissible
negation-sequences, and each side determines the other without further data.
\end{proof}

\begin{remark}[Why complementation closes the regress]
\label{rem:regress}
A positive selector is a map ``choose an element'' that must itself be chosen; in
a whole with no distinguished element the choice has no ground and regresses.
Complementation is the unique operation that fixes a part using only the whole
and the part's own absence---data already present once the whole is given. This
is why negation, and not selection, is the primitive of individuation here.
\end{remark}

\subsection{Non-instantaneity}

We now show that individuation cannot be punctual. The intuition is exact: a
separation that takes no extent leaves the part indistinguishable from the
whole, hence is no separation at all. We formalise an individuation as a finite
sequence of refinements and prove that a zero-length sequence individuates
nothing.

\begin{definition}[Separation process]
\label{def:process}
A \emph{separation process} for a target region $A$ is a finite chain of
realisable partitions
$\Part_0\succeq\Part_1\succeq\cdots\succeq\Part_k$
together with, at each stage $j$, a $\Part_j$-realisable approximant $A_j$
(a union of cells of $\Part_j$), such that $A_0=\Space$ \textup{(}the
undifferentiated whole\textup{)} and the chain is \emph{committing}: each step
$j\to j+1$ moves at least one cell out of the approximant's separator, i.e.\
$\thick_{\Part_{j+1}}(A_{j+1})<\thick_{\Part_j}(A_j)$ or
$A_{j+1}\neq A_j$. The \emph{length} of the process is $k$, the number of
committed steps.
\end{definition}

\begin{theorem}[Non-instantaneity]
\label{thm:noninstant}
Under \Cref{ax:brs} with $\Space$ connected, no separation process of length
$k=0$ individuates a part: if $k=0$ then $A_0=\Space$ and no region $A$ with
$\mu(A)>0$ and $\mu(\Space\setminus A)>0$ has been determined. Consequently every
individuation of a genuine part has length $k\ge1$, and at every stage prior to
completion a non-empty separator of measure $\ge\thick$ remains uncommitted.
\end{theorem}

\begin{proof}
At length $k=0$ the only datum is $A_0=\Space$. Then $\Space\setminus A_0=\varnothing$
has measure zero, so $A_0$ is not a region in the sense of \Cref{def:region}: it
fails $\mu(\Space\setminus A_0)>0$. No part distinct from the whole has been
determined; $A_0$ coincides with $\Space$. Hence a genuine part requires at least
one committing step, $k\ge1$.

For the residue claim, consider any stage $j<k$ with approximant $A_j$. By
\Cref{thm:negation-unique} the still-pending determination of $A$ is the
withholding of $\Neg(A)$, not yet complete at stage $j$, so $A_j\neq A$ and the
realisable separator $\Sigma_{\Part_j}(A_j)$ is non-empty. By
\Cref{thm:thickness}, $\mu(\Sigma_{\Part_j}(A_j))\ge\thick>0$. Thus at every
stage before completion a separator of measure at least the floor remains
uncommitted---the untouched residue of the individuation in progress.
\end{proof}

\begin{remark}[The residue]
\label{rem:residue-origin}
\Cref{thm:noninstant} is the generative fact behind everything that follows. An
individuation is not the flash assignment of a boundary but a process that
withholds $\Neg(A)$ step by step; because it has positive length and the whole
against which it withholds is itself only ever finitely resolved, what is left
uncommitted at the close is never empty. We call this remainder the
\emph{residue} of the individuation. \Cref{sec:residue} shows the residue cannot
be removed even in principle---not because it is small, but because removing it
would require verifying $A$ against the very totality that constitutes it.
\end{remark}

\begin{corollary}[Individuation is monotone in commitments]
\label{cor:monotone-commit}
Along any separation process the number of committed steps is non-decreasing, and
a step once committed---a cell once removed from the separator---cannot be
un-committed without itself being a further step. Hence the committed-count is a
monotone record of the process.
\end{corollary}

\begin{proof}
By \Cref{def:process} each step is committing and the chain of partitions is a
refinement chain $\Part_j\succeq\Part_{j+1}$, which is a one-directional order:
re-coarsening to undo a commitment is a distinct act, not an inversion of the
order on the original chain. Thus the count of committed steps never decreases
along the process; it is taken up again as the central object of
\Cref{sec:inquiry}.
\end{proof}
